% GENERATED FILE. Edit canonical lessons, then run npm run generate:books.
% canonical-source-hash: fnv1a64:ce15f1314006c5f7
% canonical-lessons: TA-C56-a-little, TA-C56-a-lot, TA-C56-less, TA-C56-no-need, TA-C56-let-it-be

\chapter{How Much}
\label{ch:ta-how-much}
\hlchaptermodality{56}

\begin{chapteropening}
\textbf{By the end of this chapter:} \emph{I can say how much I want in Tamil, from a little to a lot, and turn an offer down politely.}

Complete TA-C56-let-it-be and say all five words in order.
\end{chapteropening}

\section[siṟitu]{\ta{சிறிது} (siṟitu) — a little}

\label{lesson:TA-C56-a-little}



\begin{quote}
Before the new run of words: what did \emph{nakam} mean, and what did \emph{elumbu} mean?
\end{quote}

\subsection*{You'll want to know: \ta{சிறிது}}
\textbf{\ta{சிறிது}} (\emph{siṟitu}) — \textquotedblleft{}a little\textquotedblright{}.

You already have \emph{siṟiya}, \textquotedblleft{}small\textquotedblright{}. \ta{சிறிது} is the same root turned into an amount: not a small thing but a small quantity of something. Held a dish and asked how much, you say \emph{siṟitu} and you get a little.

Tamil has a \emph{koñcam} as well, which is the one you will hear across a table, and which Telugu shares as \emph{koñcem}. \ta{சிறிது} is the more formal of the two, and the one you would write down.

\begin{grammarlens}[title={What you've built}]
The first of five answers to \textquotedblleft{}how much?\textquotedblright{}.
\end{grammarlens}

\subsection*{Your turn}
Say these aloud:

\begin{itemize}
\raggedright
  \item \emph{siṟitu}
  \item it once more, slowly
  \item \emph{elumbu}, then \emph{siṟitu}, and keep the second one short
\end{itemize}

\begin{tcolorbox}[breakable,colback=teal!4,colframe=teal!35!black,title={Before you move on}]
What does \ta{சிறிது} mean? (\textquotedblleft{}a little\textquotedblright{}.) Say it once more.
\end{tcolorbox}
\section[niṟaiya]{\ta{நிறைய} (niṟaiya) — a lot}

\label{lesson:TA-C56-a-lot}



\begin{quote}
Before the new one: what did \emph{siṟitu} mean?
\end{quote}

\subsection*{You'll want to know: \ta{நிறைய}}
\textbf{\ta{நிறைய}} (\emph{niṟaiya}) — \textquotedblleft{}a lot\textquotedblright{}.

The other end of the same question. \ta{நிறைய} is built on \emph{niṟai}, to fill, so a lot in Tamil is a fullness: as much as the thing will hold and no argument about it.

Telugu says \te{ఎక్కువ} (\emph{ekkuva}) and builds it out of nothing to do with filling — a reminder that two languages can land on the same everyday answer by two quite different routes.

\begin{grammarlens}[title={What you've built}]
Two, and they are opposite ends of one question.
\end{grammarlens}

\subsection*{Your turn}
Say these aloud:

\begin{itemize}
\raggedright
  \item \emph{niṟaiya}
  \item it once more, slowly
  \item \emph{siṟitu}, then \emph{niṟaiya}, the little and the lot
\end{itemize}

\begin{tcolorbox}[breakable,colback=teal!4,colframe=teal!35!black,title={Before you move on}]
What does \ta{நிறைய} mean? (\textquotedblleft{}a lot\textquotedblright{}.) Say it once more.
\end{tcolorbox}
\section[kuṟaivu]{\ta{குறைவு} (kuṟaivu) — a shortage, less than there should be}

\label{lesson:TA-C56-less}



\begin{quote}
Before the new one: what did \emph{niṟaiya} mean?
\end{quote}

\subsection*{You'll want to know: \ta{குறைவு}}
\textbf{\ta{குறைவு}} (\emph{kuṟaivu}) — \textquotedblleft{}a shortage, less than there should be\textquotedblright{}.

\emph{kuṟai} is to fall short, and \ta{குறைவு} is the falling-short itself. That makes it a narrower word than the English \textquotedblleft{}less\textquotedblright{}: it does not merely compare two amounts, it says one of them did not come up to the mark. There is a small complaint folded inside it.

Telugu's word for the same ground is \te{తక్కువ} (\emph{takkuva}).

\begin{grammarlens}[title={What you've built}]
Three: a little, a lot, and not as much as there ought to be.
\end{grammarlens}

\subsection*{Your turn}
Say these aloud:

\begin{itemize}
\raggedright
  \item \emph{kuṟaivu}
  \item it once more, slowly
  \item \emph{niṟaiya}, then \emph{kuṟaivu}, plenty and then not enough
\end{itemize}

\begin{tcolorbox}[breakable,colback=teal!4,colframe=teal!35!black,title={Before you move on}]
What does \ta{குறைவு} mean? (\textquotedblleft{}a shortage, less than there should be\textquotedblright{}.) Say it once more.
\end{tcolorbox}
\section[tēvaiyillai]{\ta{தேவையில்லை} (tēvaiyillai) — there is no need}

\label{lesson:TA-C56-no-need}



\begin{quote}
Before the new one: what did \emph{kuṟaivu} mean?
\end{quote}

\subsection*{You'll want to know: \ta{தேவையில்லை}}
\textbf{\ta{தேவையில்லை}} (\emph{tēvaiyillai}) — \textquotedblleft{}there is no need\textquotedblright{}.

Two pieces, and the second one has been yours since the early writing lessons: \emph{tēvai} is a need, \emph{illai} is the no you learned to put on paper. Together they are how Tamil turns something down without turning the person down.

A flat \emph{illai} refuses the thing. \ta{தேவையில்லை} says there is no call for it, which leaves everybody their footing. Telugu does the same work in one shorter piece, \te{వద్దు} (\emph{vaddu}).

\begin{grammarlens}[title={What you've built}]
Four, and this one is the polite way out.
\end{grammarlens}

\subsection*{Your turn}
Say these aloud:

\begin{itemize}
\raggedright
  \item \emph{tēvaiyillai}
  \item it once more, slowly
  \item \emph{kuṟaivu}, then \emph{tēvaiyillai}, and let the whole of the second run as one word
\end{itemize}

\begin{tcolorbox}[breakable,colback=teal!4,colframe=teal!35!black,title={Before you move on}]
What does \ta{தேவையில்லை} mean? (\textquotedblleft{}there is no need\textquotedblright{}.) Say it once more.
\end{tcolorbox}
\section[ākaṭṭum]{\ta{ஆகட்டும்} (ākaṭṭum) — let it be so}

\label{lesson:TA-C56-let-it-be}



\begin{quote}
Before the new one: what did \emph{tēvaiyillai} mean?
\end{quote}

\subsection*{You'll want to know: \ta{ஆகட்டும்}}
\textbf{\ta{ஆகட்டும்}} (\emph{ākaṭṭum}) — \textquotedblleft{}let it be so\textquotedblright{}.

\emph{āku} is to become, to come about. \ta{ஆகட்டும்} is \textquotedblleft{}let it come about\textquotedblright{} — what you say when somebody proposes something and you are content to let it happen. It agrees without promising to do anything.

The \emph{ām} you have been saying for yes is that same \emph{āku} worn down by a great deal of use, which is why the two sit so alike in the mouth. Telugu's nearest of these is \te{అలాగే} (\emph{alāgē}), \textquotedblleft{}just so, alright then\textquotedblright{}.

\begin{grammarlens}[title={What you've built}]
Five answers to an offer: a little, a lot, not enough, no need, and yes — let it happen.
\end{grammarlens}

\subsection*{Your turn}
Say these aloud:

\begin{itemize}
\raggedright
  \item \emph{ākaṭṭum}
  \item it once more, slowly
  \item all five in order — \emph{siṟitu}, \emph{niṟaiya}, \emph{kuṟaivu}, \emph{tēvaiyillai}, \emph{ākaṭṭum}
\end{itemize}

\begin{tcolorbox}[breakable,colback=teal!4,colframe=teal!35!black,title={Before you move on}]
What does \ta{ஆகட்டும்} mean? (\textquotedblleft{}let it be so\textquotedblright{}.) Say it once more.
\end{tcolorbox}
